\scanpage{169}
\begin{center}- 44 -\end{center}\hfill IX
\figure{\begin{tikzpicture}\draw[line width=1pt] (0,0)--(.3,0);\end{tikzpicture}}

Pour le vérifier, il suffit de voir que l'application ensembliste sous-jacente aux morphismes $\operatorname{Spec}(k)\longrightarrow A$ représentés par les deux membres sont les mêmes, ce qui nous réduit à prouver l'égalité déduite de (4.3.5) en y composant le premier membre à gauche avec $f_q^A$, et le deuxième membre à droite avec $f^{-1}$ (puisque $f_q^A$ et $f^{-1}$ induisent l'identité sur les ensembles sous-jacents aux schémas source). Mais la formule en question est alors conséquence immédiate de (4.3.3) et de (4.3.4).

Appliquant (4.3.5) aux $u\in{}_{\ell^\nu}A(k)$, et passant à la limite sur $\nu$, on trouve que $g_\ell$ peut aussi s'interpréter comme l'endomorphisme $T_\ell(F)$ de $T_\ell(A(k))$ associé à l'endomorphisme $F$ de $A$. On sait d'autre part que $T_\ell(F)$ a un polynôme caractéristique qui est à coefficients entiers indépendants de $\ell$, en se ramenant par dévissage au cas où $A$ est, soit un schéma abélien, soit un tore, où le résultat en question était déjà utilisé dans la partie a) de notre démonstration. De plus, cette démonstration prouve également les assertions faites sur les valeurs propres, compte tenu, dans le cas d'un schéma abélien, des classiques résultats de WEIL [15].

Cela prouve 4.3.2 et achève la démonstration de 4.3, et prouve en même temps :

\uline{Corollaire} 4.4. \uline{Plaçons nous sous les conditions de} 4.3 b). \uline{Alors pour tout élément} $g$ \uline{dans} $\pi'$ \uline{dont l'image dans} $\pi_{\unsure{o|0}}$ \uline{est de la forme} $f_q^n$ \uline{avec} $n\geq0$, \uline{les valeurs propres des automorphismes induits de} $U^{\mathrm{et}}$, $U^{\mathrm{ef}}/U^{\mathrm{et}}$, $U/U^{\mathrm{ef}}$ \uline{sont des entiers algébriques dont les valeurs absolues sont respectivement} $q^n$, $q^{n/2}$ \uline{et} $1$. \uline{En ce qui concerne les deux morceaux extrêmes}, \uline{on peut dire}, \uline{plus précisément}, \uline{que les valeurs propres}
\margin{362}
\margin{362}
\margin{168}
